\section{Техническое задание}
\subsection{Основание для разработки}

Полное наименование системы:  «Программно-информационная система регистрации данных пользователей». В качестве примера реализации дипломным руководителем предложена CRM-система учёта пользователей (и не только) на примере автосервиса: веб-приложение для регистрации и учёта заявок на обслуживание, клиентов, автомобилей и пользователей с ролями, регламентированным жизненным циклом заявок и журналом изменений. Краткое наименование используемого варианта — CRM автосервиса (AutoService CRM), при этом тема ВКР формулируется в общем виде — ПИС регистрации данных пользователей.


\subsection{Цель и назначение разработки}

Цель проекта: разработать программно-информационную систему регистрации данных пользователей в виде веб-приложения (CRM), обеспечивающую учёт пользователей с ролями, регистрацию и учёт связанных сущностей (в примере реализации — заявки на обслуживание, клиенты, автомобили), регламентированный жизненный цикл операций (workflow заявок) и разграничение доступа на уровне приложения и базы данных. В качестве наглядного примера реализации выбрана CRM учёта на примере автосервиса.

С системой работают следующие группы пользователей:

\begin{itemize}
\item Администратор — полный доступ: управление пользователями и назначение ролей, заявки, клиенты, автомобили, журнал изменений. Администратор обеспечивает регистрацию и актуализацию данных пользователей системы;
\item Менеджер — создание и редактирование заявок, ведение клиентов и автомобилей, назначение мастеров, просмотр журнала изменений. В зоне ответственности менеджера — регистрация заявок и контактных данных;
\item Мастер (техник) — просмотр списка заявок, взятие заявок в работу, смена статуса назначенных ему заявок. Регистрация данных со стороны мастера сводится к фиксации смены статуса и привязки заявки к себе при «взятии в работу»
\end{itemize}

Задачами данной разработки являются:
\begin{itemize}
\item Обеспечить аутентификацию и авторизацию пользователей с явным определением ролей (admin, manager, technician);
\item Реализовать регистрацию и учёт данных: CRUD по заявкам, клиентам и автомобилям с привязкой клиент → автомобили → заявки и привязкой заявок к менеджерам и мастерам;
\item Реализовать регламентированный workflow заявок (черновик → отправлена → назначена → в работе → завершена и др.) с проверкой допустимых переходов и прав по ролям;
\item Обеспечить назначение мастера на заявку и возможность мастера «взять заявку в работу»;
\item Вести автоматический журнал изменений (аудит) по ключевым сущностям — заявки, клиенты, автомобили, пользователи;
\item Обеспечить разграничение доступа на уровне приложения (маршруты, UI) и на уровне БД (RLS), чтобы регистрация и просмотр данных строго соответствовали правам пользователя.
\end{itemize}

\subsection{Актуальность темы разработки}

Актуальность темы «Программно-информационная система регистрации данных пользователей» обусловлена повсеместной потребностью организаций в системах, которые не только хранят справочники и заявки, но и однозначно фиксируют, кто из пользователей какую операцию выполнил и в каком состоянии находятся данные. Это важно для ответственности, аудита, соблюдения требований к персональным данным и для прозрачности процессов. В малом и среднем бизнесе (в том числе в сервисных и автосервисных организациях) востребованы недорогие и понятные системы учёта без избыточной функциональности крупных корпоративных CRM. Автоматизация регистрации данных пользователей снижает потери заявок, ускоряет назначение исполнителей и даёт руководству актуальную картину. Использование современных веб-технологий (Vue 3, TypeScript, Supabase) и явной доменной архитектуры позволяет получить систему, пригодную для защиты в качестве дипломного проекта и для последующего расширения или переноса примера на другую отрасль.

\subsection{Требования к программной системе}
\subsubsection{Требования к данным}

Входные данные:  данные пользователей системы (учётные записи, профили, роли), клиентов (ФИО, телефон, email), автомобилей (марка, модель, год, госномер, привязка к клиенту), заявок (клиент, автомобиль, описание, заметки, статус, менеджер, мастер). Все вводимые и изменяемые данные должны проходить через интерфейс с проверкой прав доступа; изменения ключевых сущностей фиксируются в журнале аудита.

Выходные данные и хранимые сущности:
\begin{itemize}
	\item Список и карточки заявок с фильтрацией и поиском, с отображением связанных клиентов, автомобилей и ответственных пользователей;
	\item Справочники клиентов и автомобилей с возможностью CRUD в соответствии с ролью;
	\item Справочники клиентов и автомобилей с возможностью CRUD в соответствии с ролью;
	\item Журнал изменений (аудит) с фильтрами по сущности, действию, дате, пользователю;
	\item Текущий пользователь и его роль для отображения доступных разделов и действий в интерфейсе.
\end{itemize}

Данные хранятся в реляционной БД (PostgreSQL в Supabase). Целостность обеспечивается схемой БД и внешними ключами; безопасность — RLS-политиками и бизнес-логикой в приложении, так что пользователь не может прочитать или изменить данные, на которые у него нет прав.

\subsubsection{Функциональные требования}

\begin{itemize}
	\item Аутентификация: вход по email и паролю (Supabase Auth), выход, проверка сессии при навигации. После входа система определяет роль пользователя из профиля и ограничивает доступ в соответствии с ней;
	\item Авторизация: определение роли (admin, manager, technician), ограничение доступа к маршрутам (guards) и к действиям в UI (скрытие кнопок и разделов), проверка прав на уровне API и БД (RLS);
	\item Заявки: создание (в статусе «черновик»), просмотр списка и карточки, редактирование полей, смена статуса по регламенту, назначение мастера, «взять в работу» (для мастера), удаление (с учётом прав). Фильтрация и поиск по списку (в том числе по данным клиента и автомобиля);
	\item Клиенты: просмотр списка, создание, редактирование, удаление (доступно ролям admin, manager);
	\item Автомобили: просмотр списка, создание, редактирование, удаление, привязка к клиенту (admin, manager);
	\item Пользователи: просмотр списка пользователей системы и назначение ролей (только admin)
	\item Журнал изменений: просмотр записей аудита с фильтрами по типу сущности, действию, периоду (admin, manager)
	\item Dashboard: сводная информация и быстрые действия в зависимости от роли пользователя
\end{itemize}

\subsubsection{Моделирование вариантов использования}

На основании анализа требований в приложении реализованы следующие прецеденты.

\begin{enumerate}
	\item Авторизация — вход в систему по email и паролю; при успехе пользователь перенаправляется в приложение, роль определяется из профиля и определяет доступ к разделам;
	\item Просмотр и управление заявками — просмотр списка заявок (с учётом прав: мастер видит свои и/или общий список в зависимости от политик), открытие карточки, создание, редактирование, смена статуса, назначение мастера, взятие в работу (мастер);
	\item Управление клиентами — просмотр списка, создание, редактирование, удаление клиентов (роли admin, manager);
	\item Управление автомобилями — просмотр списка, создание, редактирование, удаление автомобилей, привязка к клиенту (admin, manager);
	\item Управление пользователями — просмотр пользователей системы, назначение ролей (только администратор);
	\item Просмотр журнала изменений — просмотр записей аудита с фильтрами (admin, manager);
	\item Выход — завершение сессии и перенаправление на страницу входа
\end{enumerate}

На рисунках 2.1 и 2.2 представлены диаграммы прецедентов

\begin{figure}[ht]
	\includegraphics[width=1\linewidth]{user-use-case}
	\caption{Диаграмма прецедентов пользователя (менеджер / мастер) для программы}
	\label{user-use-case:image}
\end{figure}

\begin{figure}[ht]
	\includegraphics[width=1\linewidth]{admin-use-case}
	\caption{Диаграмма прецедентов администратора для программы}
	\label{admin-use-case:image}
\end{figure}

\textbf{Сценарий прецедента «Авторизация»:}

\begin{itemize}
	\item Основной исполнитель: пользователь (сотрудник организации);
	\item Предусловие: приложение развёрнуто, Supabase настроен, учётная запись пользователя создана;
	\item Описание: пользователь вводит email и пароль, нажимает «Войти». Система проверяет учётные данные через Supabase Auth, загружает профиль с ролью (admin, manager, technician). При успехе выполняется переход в приложение (например, на Dashboard), меню и маршруты формируются с учётом роли; при отказе отображается сообщение об ошибке.
	\item Результат: пользователь авторизован, роль определена, доступ к разделам ограничен по роли в соответствии с требованиями регистрации данных пользователей.
\end{itemize}

\textbf{Сценарий прецедента «Смена статуса заявки»:}

\begin{itemize}
	\item Основной исполнитель: менеджер или мастер (в зависимости от перехода и назначения);
	\item Предусловие: пользователь авторизован, заявка открыта, переход из текущего статуса в целевой допустим по регламенту (workflow) и по роли (например, завершать может менеджер или назначенный мастер);
	\item Описание: пользователь выбирает новый статус (например, «Завершена») и при необходимости вводит заметку. Система проверяет допустимость перехода по графу статусов и права текущего пользователя. Выполняется вызов API смены статуса, обновление данных в БД, запись в журнал аудита; обновлённые данные отображаются в интерфейсе; 
	\item Результат: статус заявки изменён, в журнал аудита записано изменение с указанием пользователя и времени. 
\end{itemize}

\textbf{Бизнес-правила:}

\begin{itemize}
	\item Переходы между статусами заявок заданы явно (конечный автомат, state machine); недопустимые переходы блокируются в UI и на уровне API/БД;
	\item Назначение мастера допустимо для заявок в статусе «Отправлена» или «Назначена». Взять заявку в работу может только мастер; при этом заявка переводится в статус «Назначена» с заполнением поля исполнителя (technician\_id) текущим пользователем (или в «В работе» в зависимости от сценария);
	\item Удаление и массовые операции с заявками доступны только при наличии соответствующих прав (администратор, менеджер);
	\item Все изменения ключевых сущностей (заявки, клиенты, автомобили, пользователи — в том числе смена роли) фиксируются в журнале изменений (аудит).
\end{itemize}

\subsubsection{Требования к безопасности}

\textbf{Аутентификация:} через Supabase Auth; пароли не хранятся в приложении, сессии управляются сервисом.\linebreak
\textbf{Авторизация по ролям:} проверка в маршрутах (guards), в UI (скрытие кнопок и разделов) и на уровне БД (RLS). RLS является последним рубежом: пользователь не может прочитать или изменить строки, запрещённые политиками, даже при прямом обращении к API.\linebreak
\textbf{Защита от несанкционированного доступа:} только авторизованные пользователи получают доступ к данным; уровень доступа определяется ролью (admin, manager, technician) и при необходимости привязкой к записи (например, мастер видит только свои заявки по полю technician\_id в зависимости от настроек политик).\linebreak
\textbf{Обработка ошибок:} программа должна корректно обрабатывать ошибки (сеть, отказ API, неверные данные) и не раскрывать внутренние детали (стек вызовов, структура БД) в сообщениях конечному пользователю.

\subsubsection{Требования к программному обеспечению}

Для разработки и эксплуатации используются:

\begin{itemize}
	\item \textbf{Frontend:} Node.js 18+, npm или yarn; Vue 3 (Composition API), TypeScript, Vite, Pinia, Vue Router, PrimeVue (или выбранная UI-библиотека);
	\item \textbf{Backend / BaaS:} проект Supabase (PostgreSQL, Auth, REST API, RLS). Доступ по HTTPS;
	\item \textbf{Среда разработки:} едактор кода (например, VS Code), браузер с поддержкой современного JavaScript (ES2020+);
	\item \textbf{Сборка:} сборка статического SPA (npm run build); размещение на любом хостинге статики (Vercel, Netlify и т.д.) или на сервере организации;
	\item Поддержка русской раскладки клавиатуры и отображение интерфейса на русском языке.
\end{itemize}

\subsubsection{Требования к аппаратному обеспечению}

Для работы приложения достаточно: ПК или мобильное устройство с браузером, поддерживающим современный JavaScript; подключение к интернету для доступа к Supabase. Рекомендуется не менее 2 ГБ ОЗУ, разрешение экрана не ниже 1024×768 для комфортной работы с таблицами и формами.

\subsubsection{Требования к оформлению документации}

Требования к стадиям разработки и программной документации устанавливаются ГОСТ 19.102–77 и внутренними стандартами вуза. Документация проекта включает описание архитектуры, руководство по развёртыванию и при необходимости руководство пользователя.
